% -----------------------------------------------------------------------------
%  entries/q030.tex
%
%  One file per problem, named after its permanent ID.  The ID never changes
%  and never gets reused, even if the entry later moves to a different day.
%
%  Keys are comma-separated, so any value containing a comma must be braced:
%      topic = {Trace norm, Hilbert-Schmidt norm}
% -----------------------------------------------------------------------------
\begin{entry}{Q030}{Charlie's help}

\entrymeta{
  topic   = {Multipartite entanglement, measurement on a third party},
  est     = {$\sim$30 minutes (including write-up)},
  origin  = Original,          % Sourced / Extracted / Adapted / Original
  source  = Original problem,
  prereq  = Schmidt decomposition,
  status  = drafted,           % drafted / solved / written-up / published
}

\begin{problem}
Tracing out one qubit of $|\mathrm{GHZ}\rangle$ leaves the other two in a
separable state, whereas the two-qubit marginal of $|W\rangle$ is entangled;
in this sense $W$ is the more robust of the two under the loss of a particle
(D\"ur, Vidal, Cirac, Phys.\ Rev.\ A \textbf{62}, 062314 (2000)). Losing a
qubit, however, is not the only thing that can happen to it: its holder may
instead measure it and announce the result. This problem asks which of the
two states is more useful when the third party cooperates in this way.

Three qubits $A$, $B$, $C$, held by Alice, Bob and Charlie, are in one of
the states
\begin{equation}\label{eq:q030-states}
  |\mathrm{GHZ}\rangle = \tfrac{1}{\sqrt2}\big(|000\rangle + |111\rangle\big),
  \qquad
  |W\rangle = \tfrac{1}{\sqrt3}\big(|001\rangle + |010\rangle + |100\rangle\big),
\end{equation}
with the qubits ordered $ABC$. Charlie measures $C$ in an orthonormal basis
$\{|\varphi_0\rangle, |\varphi_1\rangle\}$ of $\mathbb{C}^2$ of his choice
and announces the outcome $k$. For a three-qubit state $|\psi\rangle$,
outcome $k$ occurs with probability
$p_k = \big\|(\mathds{1}\otimes\mathds{1}\otimes\langle\varphi_k|)
|\psi\rangle\big\|^2$ and, when $p_k > 0$, leaves $AB$ in the pure state
\begin{equation}\label{eq:q030-post}
  |\psi_k\rangle = \frac{1}{\sqrt{p_k}}\,
  (\mathds{1}\otimes\mathds{1}\otimes\langle\varphi_k|)\,|\psi\rangle .
\end{equation}
A normalized two-qubit vector is \emph{maximally entangled} if its reduced
state on $A$ is $\mathds{1}/2$. You may use the following fact: writing a
normalized $v = \sum_{j,k} M_{jk}|jk\rangle \in \mathbb{C}^2\otimes
\mathbb{C}^2$, its reduced state on $A$ is $MM^\dagger$, so
\begin{equation}\label{eq:q030-max}
  v \text{ is maximally entangled}
  \quad\Longleftrightarrow\quad
  MM^\dagger = \tfrac12\,\mathds{1}.
\end{equation}

\begin{enumerate}
  \item[(a)] Find a basis for Charlie such that, for $|\mathrm{GHZ}\rangle$,
  every outcome leaves Alice and Bob maximally entangled.

  \item[(b)] Show that for $|W\rangle$ an outcome leaves Alice and Bob
  maximally entangled only if the corresponding basis vector is $|0\rangle$
  up to a phase. Conclude that no basis has the property of (a), and compute
  the probability of the maximally entangled outcome when it occurs.
\end{enumerate}
\end{problem}

\begin{solution}[a)]
One can use the fact that the Bell states $\ket{\Phi^{\pm }} = \frac{1}{\sqrt{2}} \left(\ket{00} \pm \ket{11}\right)$ are maximally entangled and notice that, for the $X-$Pauli basis $\ket{\pm} = \frac{1}{\sqrt{2}} (\ket{0}\pm \ket{1})$, we have 
\begin{equation}
    (\mathds{1}\otimes\mathds{1}\otimes\bra{\pm}) \ket{GHZ} = \frac{1}{\sqrt{2}} \ket{\Phi^{\pm}}
\end{equation}

Thus, clearly being one possible choice of basis that always leaves the outcome maximally entangled

\end{solution}

\begin{solution}[b)]
Consider the most general way we can parametrize an orthonormal basis on the one-qubit space is: 
\begin{equation}
    \ket{\varphi_0} = \cos(\theta) \ket{0} +e^{i\phi} \sin(\theta) \ket{1}, \quad  \ket{\varphi_1} = \sin(\theta) \ket{0} - e^{i\phi} \cos(\theta) \ket{1}
\end{equation}

Then the projections are
\begin{equation}
\begin{split}
    \bra{\phi_0} \ket{W} = \frac{1}{\sqrt{3}} \big[ \cos(\theta)\ket{01} + \cos(\theta)\ket{10} + e^{-i\phi}\sin(\theta)\ket{00} \big] \\
    \bra{\phi_1} \ket{W} = \frac{1}{\sqrt{3}} \big[ \sin(\theta)\ket{01} + \sin(\theta)\ket{10} - e^{-i\phi}\cos(\theta)\ket{00} \big]
\end{split}
\end{equation}
which leads to probabilities
\begin{equation}
    p_0 = \frac{1}{3} (\cos(\theta)^{2} +1), \quad p_1= \frac{1}{3} (\sin(\theta)^{2} +1)
\end{equation}
such that the corresponding matrices ( where $M^{(k)}$ is associated with the state $\ket{\psi_k}$) are:
\begin{equation}
\begin{split}
    M^{(0)} = \frac{1}{\sqrt{\cos(\theta)^{2} +1}}\begin{pmatrix}
        e^{-i\phi}\sin(\theta) & \cos(\theta) \\
        \cos(\theta) & 0
    \end{pmatrix} \\
    M^{(1)} = \frac{1}{\sqrt{\sin(\theta)^{2} +1}}\begin{pmatrix}
        e^{-i\phi}\cos(\theta) & \sin(\theta) \\
        \sin(\theta) & 0
    \end{pmatrix}
\end{split}
\end{equation}


Now we can check whether the condition \eqref{eq:q030-max} can be satisfied simultaneously for both matrices by calculating $MM^{\dagger}$ for both of them.

From their expression
\begin{equation}
\begin{split}
    M^{(0)}(M^{(0)})^{\dagger} = \frac{1}{\cos(\theta)^{2} +1}\begin{pmatrix}
        1 & \e^{-i\phi}\sin(\theta)\cos(\theta) \\
        \e^{i\phi}\sin(\theta)\cos(\theta) & \cos(\theta)^2
    \end{pmatrix} \\
    M^{(1)}(M^{(1)})^{\dagger} = \frac{1}{\sin(\theta)^{2} +1}\begin{pmatrix}
        1 & \e^{-i\phi}\sin(\theta)\cos(\theta) \\
        \e^{i\phi}\sin(\theta)\cos(\theta) & \sin(\theta)^2
    \end{pmatrix}.
\end{split}
\label{q030: conjugate matrices}
\end{equation}

But this implies that we cannot have both resulting matrices proportional to the identity simultaneously. For the first case, we would need $\theta = n\pi$ ($\sin{\theta}=0$), whereas for the second we would need $\theta = n\pi +\pi/2$ ($\cos{\theta}=0$). 

In both cases, this implies that we only achieved a maximally entangled state by projecting through $\bra{0}$ up to a phase. At the same time, the probability is $\frac{2}{3}$. 


\end{solution}

% \begin{remark}
% \end{remark}

\begin{takeaway}
\end{takeaway}

%  Everything below is optional.  An entry with no attempt block is complete.
%  Record only what would change what you do next time: typos and arithmetic
%  slips you caught yourself are artefacts of the medium, not facts about the
%  mathematics.  Add "public," to show this block in the public build.
% \attempt{
%   date       = ,               % ISO, YYYY-MM-DD -- required if the block is used
%   time       = ,
%   difficulty = ,               % 1-5, perceived
%   howitwent  = ,
%   insight    = ,
%   note       = ,
% }

\end{entry}